\documentclass[12pt]{article}
\usepackage{amsmath, amssymb, amsthm}
\usepackage{parskip}
\usepackage{geometry}
\usepackage{enumitem}
\usepackage{fancyhdr}

\geometry{letterpaper, margin=1in}
\setlength{\headheight}{15pt}

% Header setup
\pagestyle{fancy}
\fancyhf{}
\lhead{Math 311 Fall 2026}
\rhead{Jeffrey Lu}
\cfoot{\thepage}

\title{Homework Assignment 4 \\ \large Math 311, Fall 2026 \\ Rutgers University}
\author{}
\date{}

\begin{document}

\maketitle

\textbf{Exercise 2.2.7}
\begin{enumerate}[label=(\alph*)]
    \item The sequence \((-1)^n\) is only frequently in the set \(\{1\}\)
          since for all \(N \in \mathbb{N}\), there exists even \(n \geq N\) so that
          \((-1)^n = 1\). However, it is not eventually in \(\{1\}\) since
          we can also always choose an odd \(n \geq N\) so that \((-1)^n = -1\).
    \item Eventually is the stronger definition since eventually \(\implies \) frequently,
          however the converse is not true as shown in part (a).
    \item A sequence \((a_n)\) converges if for every \(\epsilon > 0 \), the sequence
          is eventually in \(V_{\epsilon}(a)\).
    \item Consider the sequence
          \[
              (x_n) = \begin{cases}
              2, & \text{if $n$ is even}   \\
              0, & \text{if $n$ is odd}
              \end{cases}
          \]
          Then we have an infintie number of terms equal to 2.
          Similarly to part (a), the sequence \((x_n)\) is frequently in \((1.9, 2.1)\)
          since for all \(N \in \mathbb{N}\), there exists even \(n \geq N\) so that
          \((x_n) = 2\). However, it is not eventually in \((1.9, 2.1)\) since
          we can always choose an odd \(n \geq N\) so that \((x_n) = 0\).
\end{enumerate}

\textbf{Exercise 2.3.2}
Suppose \( (x_n) \to 2 \).
\begin{enumerate}[label=(\alph*)]
    \item Let \( \epsilon > 0 \). Choose \( N \) such that, for \( n \geq N \),
          \[
              |x_n - 2| < \frac{3\epsilon }{2}.
          \]
          Then, for \( n \geq N \) we have
          \[
              \begin{split}
                  \left\vert \frac{2x_n - 1}{3} - 1 \right\vert & = \left\vert \frac{2x_n - 4}{3} \right\vert \\
                                                                & = \frac{2}{3} \vert x_n - 2 \vert           \\
                                                                & < \frac{2}{3} \cdot \frac{3\epsilon }{2}    \\
                                                                & = \epsilon .
              \end{split}
          \]
          Thus,
          \[
              \frac{2x_n - 1}{3} \to  1.
          \]

    \item Let \( \epsilon > 0 \). Choose \(  N \) such that, for \(  n \geq N \),
          \[
              \vert x_n - 2 \vert < \min\{1, 2\epsilon \}.
          \]
          Then, for \( n \geq N \) we have
          \[
              \begin{split}
                  \left\vert \frac{1}{x_n} - \frac{1}{2} \right\vert & = \frac{\vert 2 - x_n \vert }{2\vert x_n \vert } \\
                                                                     & < \frac{\vert 2 - x_n \vert }{2}                 \\
                                                                     & < \frac{2\epsilon }{2}                           \\
                                                                     & = \epsilon.
              \end{split}
          \]
          Thus,
          \[
              \frac{1}{x_n} \to \frac{1}{2}.
          \]
\end{enumerate}

\textbf{Exercise 2.3.5}

Let
\[
    (z_n) = (x_1, y_1, x_2, y_2, x_3, y_3, \dots ).
\]
That is,
\[
    z_{2n-1} = x_n, \qquad z_{2n} = y_n.
\]


Suppose first that \( (z_n) \to L\). Let \( \epsilon > 0 \) and choose \( N \in \mathbb{N} \) such that
\[
    k \geq N \implies |z_k - L| < \epsilon .
\]
For every \( n \geq N \), it follows that \( 2n - 1 \geq N \) and \( 2n \geq N \).
Hence
\[
    |x_n - L| =  \vert z_{2n-1} - L \vert < \epsilon , \qquad \vert y_n - L \vert = \vert z_{2n} - L \vert < \epsilon .
\]
Thus \( x_n \to L \) and \( y_n \to  L \).

Conversely, suppose \( x_{n} \to L \) and \( y_n \to  L \). Let \( \epsilon > 0 \) and choose
\( N_x, N_y \in \mathbb{N} \) such that
\[
    n \geq N_x \implies \vert x_n - L \vert < \epsilon , \qquad n \geq N_{y} \implies  \vert y_n - L \vert < \epsilon .
\]
Let \( M = \max \{N_x, N_{y} \} \), and \( N = 2M \). For any \( m \geq N \):

\begin{itemize}
    \item If \( k = 2n \), then \( n \geq M \), so
          \[
              \vert z_k - L \vert = \vert y_n - L \vert < \epsilon .
          \]
    \item If \( = 2n - 1 \), then \( n \geq M \), so
          \[
              \vert z_k - L \vert = \vert x_n - L \vert < \epsilon .
          \]
\end{itemize}

Therefore, \( z_n \to L \).

\textbf{Exercise 2.3.6} Take \( (\frac{1}{n}) \to 0\) as given, and consider
\[
    \begin{split}
        b_n & = n - \sqrt{n^2 + 2n}                            \\
            & = \frac{n^2 - (n^2 + 2n)}{n + \sqrt{n^{2} + 2n}} \\
            & = \frac{-2n}{n + \sqrt{n^2 + 2n}}                \\
            & = \frac{-2n}{n + n \sqrt{1 + \frac{2}{n}}}       \\
            & = \frac{-2n}{n(1 + \sqrt{1 + \frac{2}{n}})}      \\
            & = \frac{-2}{1 + \sqrt{1 + \frac{2}{n}}}.
    \end{split}
\]
Since \( (\frac{1}{n}) \to 0 \), by Algebraic Limit Theorem we have that
\[
    1 + \frac{2}{n} \to  1,
\]
and by exercise 2.3.1, we have
\[
    \sqrt{1 + \frac{2}{n}} \to  \sqrt{1} = 1.
\]
Thus the denominator converges
\[
    1 + \sqrt{1 + \frac{2}{n}} \to  1 + 1 = 2.
\]
Using the quotient rule of the Algebraic Limit Theorem,
\[
    b_n \to -\frac{2}{2} = -1.
\]

\textbf{Exercise 2.3.8}
Let \( (x_n) \to  x\) and let \( p(x) \) be a polynomial.
\begin{enumerate}[label=(\alph*)]
    \item Notice that
          \[
              p(x_n) = a_0 + a_1 x_n + a_2 x_n^2 + \dots + a_k x_n^k \quad \text{for some } k \in \mathbb{N}.
          \]
          Since \( (x_n) \to x \), the Algebraic Limit Theorem says
          \[
              a_k x_n^k \to a_k x^k \qquad \text{for every } k \in \mathbb{N}.
          \]
          Then applying the sum rule for every element in \( p(x_n) \), we have
          \[
              p(x_n) \to  a_0 + a_1 x + \dots a_k x^k = p(x).
          \]
    \item Let \( x_n = \frac{1}{n} \) and define
          \[
              f(t) = \begin{cases}
                  1, & t \neq 0 , \\
                  0, & t = 0.
              \end{cases}
          \]
          Then \( x_n \to 0 \). But since \( x_n \neq 0 \) for every \( n \), we have \( f(x_n) = 1 \), and so
          \[
              f(x_n) \to 1, \qquad \text{while} \qquad f(0) = 0.
          \]
          Thus
          \[
              f(x_n) \not \to f(x).
          \]
\end{enumerate}

\textbf{Exercise 2.3.11}
\begin{enumerate}[label=(\alph*)]
    \item Assume \( (x_n) \to L\).
          Define
          \[
              y_n = \frac{x_1 + x_2 + \dots + x_n}{n}.
          \]
          Let \( \epsilon >0 \) and choose \( N \in \mathbb{N} \) such that
          \[
              n \geq N \implies |x_n - L| < \frac{\epsilon}{2}.
          \]
          Then
          \[
              \begin{split}
                  |y_n - L| & = \frac{1}{n} \sum_{k=1}^n (x_k - L)                                                                \\
                            & \leq \frac{1}{n} \sum_{k=1}^{n} \vert x_k - L \vert                                                 \\
                            & = \frac{1}{n} \sum_{k=1}^{N-1} \vert x_k - L \vert + \frac{1}{n} \sum_{k=N}^{n} \vert x_k - L \vert \\
              \end{split}
          \]
          Define
          \[
              C = \sum_{k=1}^{N-1} \vert x_k - L \vert.
          \]
          Then
          \[
              \begin{split}
                  \vert y_n - L \vert & \leq \frac{C}{n} + \frac{1}{n} \sum_{k=N}^{n} \vert x_k - L \vert \\
                                      & \leq \frac{C}{n} + \frac{n - N + 1}{n} \cdot \frac{\epsilon}{2}   \\
                                      & \leq  \frac{C}{n} + \frac{\epsilon}{2}.
              \end{split}
          \]
          Choose \( M \geq N \) such that
          \[
              M > \frac{2C}{\epsilon}.
          \]
          Then for every \( n \geq M \),
          \[
              \frac{c}{n} < \frac{c}{\frac{2c}{\epsilon}} = \frac{\epsilon}{2}.
          \]
          So
          \[
              \vert y_n - L \vert < \frac{\epsilon}{2} + \frac{\epsilon}{2} = \epsilon.
          \]
          Therefore
          \[
              y_n \to L.
          \]
    \item Consider \( x_n = (-1)^n \), this sequence does not converge. Then
          \[
              y_n = \frac{1}{n} \sum_{k=1}^{n} (-1)^k.
          \]
          Notice that even terms are 1 while odd terms are -1. So each term will cancel each other out.
          Thus
          \[
              y_n \to 0.
          \]
\end{enumerate}


\textbf{2.4.1}
\begin{enumerate}[label=(\alph*)]
    \item Define \( x_1 = 3 \) and
          \[
              x_{n+1} = \frac{1}{4-x_n}.
          \]
          We want to show that \( 0 \leq x_n \leq 3 \) by induction.
          For the base case: \( x_1 = 3 \). Suppose \( 0 < x_n \leq 3 \).
          Then \( 4 - x_n \geq 1 \), so
          \[
              0 < x_{n+1} \leq 1 \leq 3.
          \]
          Thus \( 0 < x_n \leq 1 \leq 3 \) by induction.

          Now, we want to show that \( (x_n) \) is decreasing by induction.
          Notice that \( x_2 = \frac{1}{4-3} = 1 \).
          So \( x_2 < x_1 \) holds.
          Suppose \( x_{n+1} < x_n \). Then
          \[
              4 - x_{n+1} > 4 - x_n.
          \]
          So
          \[
              x_{n+2} = \frac{1}{4 - x_{n+1}} < \frac{1}{4 - x_n} = x_{n+1}.
          \]
          That is, \( x_{n+2} < x_{n+1} \).
          Thus \( (x_n) \) is decreasing by inuction.

          Since \( x_n \) is montone decreasing and bounded, by the Monotone Convergence, \( x_n \) converges to some \( L \).

    \item From part (a), we have that \( x_n \to L \).
          Let \( \epsilon > 0  \) and choose \( N \in \mathbb{N} \) such that
          \[
              n \geq N \implies \vert x_n - L \vert < \epsilon.
          \]
          Since \( n + 1 > n \geq N \), we also have
          \[
              |x_{n+1} - L | < \epsilon.
          \]
          Thus
          \[
              x_{n+1} \to L.
          \]
    \item Taking the limit, we have
          \[
              L = \frac{1}{4 - L},
          \]
          which simplifies to
          \[
              L^2 - 4L + 1 = 0.
          \]
          Using the quadratic formula, we get
          \[
              L = 2 \pm \sqrt{3}.
          \]
          Since \( x_n \leq 1 \) when \( n \geq 2 \), we choose
          \[
              L = 2 - \sqrt{3}.
          \]
          Thus the final limit is
          \[
              x_n \to 2 - \sqrt{3}.
          \]
\end{enumerate}

\textbf{Exercise 2.4.2}
\begin{enumerate}[label=(\alph*)]
    \item The issue is we never proved that the sequence converges, so we can't just take the limits.
    \item Let \( y_1 = 1 \) and
          \[
              y_{n+1} = 3 - \frac{1}{y_n}.
          \]
          We want to show that \( 1 \leq y_n < 3 \) for every \( n \in \mathbb{N} \), so we proceed by induction.
          First, \( y_1 = 1 \), so the base case holds.
          Suppose \( 0 \leq y_k < 3 \) for some \( k \in \mathbb{N} \).
          Then
          \[
              0 < \frac{1}{3} \leq \frac{1}{y_k} \leq 1.
          \]
          So
          \[
              2 \leq 3 - \frac{1}{y_k} < 3.
          \]
          That is,
          \[
              1 \leq y_{k+1} < 3.
          \]
          So \( (y_n) \) is bounded for every \( n \in \mathbb{N} \) by induction.

          Now we want to show \( y_n \) is increasing as well.
          First,
          \[
              y_2 = 3 - \frac{1}{1} = 2 > y_1.
          \]
          Suppose \( y_{k+1} > y_{k} \) holds for some \( k \in \mathbb{N} \).
          Then
          \[
              \frac{1}{y_{k+1}} < \frac{1}{y_k},
          \]
          and
          \[
              -\frac{1}{y_{k+1}} > -\frac{1}{y_k}.
          \]
          Thus
          \[
              y_{k+2} = 3 - \frac{1}{y_{k+1}} > 3 - \frac{1}{y_k} = y_{k+1}.
          \]
          So \( (y_n) \) is increasing for every \( n \in \mathbb{N} \) by induction.

          Therefore, by monotone convergence Theorem since \( y_n \) is bounded and monotone, \( y_n \) is convergent.
          Hence we can use part (a)'s method to find the explicit limit.
\end{enumerate}

\textbf{Exercise 2.4.3}
\begin{enumerate}[label=(\alph*)]
    \item Define \( x_1 = \sqrt{2} \) and
          \[
              x_{n+1} = \sqrt{2 + x_n}.
          \]
          We want to show \( \sqrt{2} \leq x_n < 2 \) for every \( n \in \mathbb{N} \).
          First, \( x_1 = \sqrt{2} \), so the base case holds.
          Suppose \( \sqrt{2} \leq x_k < 2 \) for some \( k \in \mathbb{N} \).
          Then
          \[
              \sqrt{2} < x_{k+1} = \sqrt{2 + x_k} < \sqrt{4} = 2.
          \]
          So \( \sqrt{2} \leq x_n < 2 \) holds for every \( n \in \mathbb{N} \) by induction.

          Next, we want to show that \( x_n \) is increasing.
          First,
          \[
              x_2 = \sqrt{2 + \sqrt{2}} > \sqrt{2} = x_1.
          \]
          So the base case holds.
          Suppose \( x_{k+1} > x_k \) for some \( k \in \mathbb{N} \).
          Then
          \[
              x_{k+2} = \sqrt{2 + x_{k+1}} > \sqrt{2 + x_k} = x_{k+1}.
          \]
          Thus the sequence is increasing by induction.

          Since \( x_n \) is increasing and bounded above, the monotone convergence theorem says \( x_n \to L \),
          where \( \sqrt{2} \leq L < 2 \). Taking the limits gives
          \[
              L = \sqrt{2 + L} \implies  \qquad (L-2)(L+1) = 0.
          \]
          Since \( L \geq  \sqrt{2} \), we get
          \[
              x_n \to 2.
          \]


    \item Define \( y_1 = \sqrt{2} \) and
          \[
              y_{n+1} = \sqrt{2y_n}.
          \]
          We want to show \( \sqrt{2} \leq y_n < 2 \) for every \( n \in \mathbb{N} \).
          First, \( y_1 = \sqrt{2} \), and so the base case holds.
          Suppose \( \sqrt{2} \leq y_k < 2 \) for some \( k \in \mathbb{N} \).
          Then
          \[
              2\sqrt{2} \leq 2y_k < 2 \implies \sqrt{2\sqrt{2}} \leq \sqrt{2y_k} < 2.
          \]
          Thus
          \[
              \sqrt{2} \leq y_{k+1} < 2.
          \]
          So \( \sqrt{2} \leq y_n < 2 \) holds for every \( n \in \mathbb{N} \) by induction.

          Now we show \( y_n \) is increasing. First,
          \[
              y_2 = \sqrt{2\sqrt{2}} > \sqrt{2} = y_1,
          \]
          and so the base case holds.
          Suppose \( y_{k+1} > y_k \) for some \( k \in \mathbb{N} \).
          Then
          \[
              y_{k+2} = \sqrt{2y_{k+1}} > \sqrt{2y_k} = y_{k+1}.
          \]
          Thus \( y_n \) is increasing by induction.

          Since \( y_n \) is increasing and bounded above by \( 2 \), the monotone convergence theorem says
          \( y_n \to L \).
          Taking limits,
          \[
              L = \sqrt{2L} \implies L(L-2) = 0.
          \]
          Since \( L > \sqrt{2} \), the limit is
          \[
              y_n \to 2.
          \]
\end{enumerate}

\textbf{Exercise 2.4.8}
\begin{enumerate}[label=(\alph*)]
    \item
\end{enumerate}





















\end{document}